\section{Related work}\label{sec:related}
\paragraph{IBLT and set reconciliation.} IBLT\cite{gm} and the invertible Bloom filter\cite{eppstein2011whatsdifference} list set differences by peeling. Decoding success is determined by the 2-core of a random hypergraph\cite{molloy,dw,dm}. At low load, duplicate edge pairs dominate failure\cite{gm,yw}, and the probability of a linear-sized core at subcritical load is exponentially small\cite{kim2008poisson,dn08}. Subsequent work improves the hypergraph structure\cite{lazaro2021irregular}, develops rateless codes\cite{yang2024rateless,lazaro2023ratecompatible}, or predicts and breaks the 2-core after decoding failure\cite{gadurek2026breaking}. Belazzougui, Kucherov, and Walzer\cite{belazzougui2024better} use finite-field markers to distinguish the direction of the difference and make the failure probability exponentially small. These works seek to reduce failure or resume decoding after failure. We read the same table after failure to determine the next-round capacity.

\paragraph{Estimating the difference size.} A common approach to estimating $d$ is to send a separate estimator, such as the strata estimator\cite{eppstein2011whatsdifference}, min-wise hashing\cite{broder2000minwise}, the odd sketch\cite{mitzenmacher2014odd}, or GXBits\cite{jia2024gxbits}. PBS\cite{gong2021pbs} first estimates $d$ using 128 independent ToW sketches, then sets capacity to a fixed multiple of the estimate. It transmits and accounts for the estimator and reconciliation structure separately. We obtain the estimate from the same counts used for decoding, avoiding a separate estimator. This reuse requires analysis conditional on decoding failure: the estimate is needed only after failure, and failure selects the mapping.

\paragraph{Random-sign sketches and second-moment estimation.} AMS\cite{alon1999spacecomplexity} estimates $F_2$ using a sum of random signs. CountSketch\cite{charikar2004countsketch} supports point queries with random-sign counters. Thorup and Zhang\cite{thorup2004tabulation} analyze its second-moment extension and give a centered second-moment estimate for a single counter array. Fast-AGMS\cite{cormode2005fastagms} uses the same structure for distributed queries. $\widehat d$ has the same algebraic form as these estimators. Its conditional unbiasedness given the mapping also follows the structure of the Hutchinson trace estimator\cite{hutchinson1990trace}. These sketches carry no decodable keySum and do not address the distribution conditional on decoding failure. Fischer and Nakos\cite{fischer2025sparse} combine hypergraph peeling and a random-sign CountSketch in a single measurement. An independent short sketch supplies the tail norm, which their analysis treats as known. We encode the shared random signs in the IBLT's count, keySum, and checksum. The same table supports peeling and, after failure, provides an estimate independent of the difference direction, together with a post-failure upper bound valid at all loads.

\paragraph{Reading the same table after failure.} Prior work\cite{wu2026selfsizing} reads count energy after first-round failure to set the second-round capacity. Its analysis models sign composition with iid signs and obtains a conditional mean equal to $d$. Its upper-bound rule uses a chi-square approximation under deep overload and per-configuration simulation calibration at other loads. We encode random signs in all three IBLT fields, making these properties hold exactly for every fixed mapping. We also derive a post-failure distribution and upper bound valid at all loads.

\paragraph{Conditional limits and discrete approximations.} For bounded $c$ (\eqref{eq:c}), we derive the limit distribution using the Chen--Stein method\cite{arratia1989chen} for the Poisson approximation of sparse overlap. The Skellam distribution\cite{skellam1946} describes the difference between two types of counts. At low load, we replace the failure-conditioned distribution with the planted distribution, which is the distribution size-biased by the duplicate edge pair count\cite{arratia2019sizebias}. When $d\gg\sqrt M$, we analyze the random-sign quadratic form given the mapping using the cumulant method on a dependency graph\cite{janson1988cumulant}. This requires only concentration bounds\cite{warnke2016typical,combes2015mcdiarmid} and degree control for the overlap graph. We prove a Gaussian limit uniform over deterministic coefficients using an edge-exposure martingale CLT\cite{hall1980martingale}. A general quadratic form CLT is given in\cite{dejong}. The upper-bound rule uses classical test inversion\cite{neyman1937}.
